% Source: https://artofproblemsolving.com/community/c3h2171804_largest_value
% Stable identity: aops-topic-2171804-q01
Factor out $x-b$:
\[(x-b)\left((x-a)+(x-c)\right)=(x-b)(2x-a-c)=0,\]
so the solutions are $x=b$ and $x=\dfrac{a+c}{2}$, with sum $b+\dfrac{a+c}{2}$. (If the two solutions coincide, this is still the sum of the roots counted with multiplicity.) To maximize it, take $b=9$ and then $a+c$ as large as possible with $a,c$ different from each other and from $9$: $a+c=8+7=15$. The maximum is
\[9+\dfrac{15}{2}=\dfrac{33}{2},\]
so $p=33$, $q=2$, and $p+q=\boxed{35}$.
